\documentclass[11pt]{article}

\usepackage[margin=1in]{geometry}
\usepackage{amsmath, amssymb, amsthm}
\usepackage{mathrsfs}
\usepackage{enumitem}
\usepackage[dvipsnames]{xcolor}
\usepackage{titlesec}
\usepackage{fancyhdr}
\usepackage[hidelinks]{hyperref}
\usepackage{tcolorbox}

\tcbuselibrary{breakable}

% Shortcuts
\newcommand{\R}{\mathbb{R}}
\newcommand{\N}{\mathbb{N}}
\newcommand{\Z}{\mathbb{Z}}
\newcommand{\Q}{\mathbb{Q}}
\newcommand{\C}{\mathbb{C}}
\newcommand{\eps}{\varepsilon}

% Theorem environments
\theoremstyle{definition}
\newtheorem*{recipe}{Strategy}

% Colors matching lecture notes
\definecolor{lectureorange}{RGB}{255, 200, 120}
\definecolor{babyblue}{RGB}{173, 216, 230}
\definecolor{lightred}{RGB}{255, 180, 180}

% Colored boxes (same scheme as notes: orange=overview, blue=section, red=notes)
\newtcolorbox{intuition}[1][]{colback=lectureorange!30, colframe=lectureorange!80, title=Intuition, fonttitle=\bfseries, breakable, #1}
\newtcolorbox{application}[1][]{colback=babyblue!30, colframe=babyblue!80, title=How to Apply, fonttitle=\bfseries, breakable, #1}
\newtcolorbox{pitfall}[1][]{colback=lightred!30, colframe=lightred!80, title=Common Pitfall, fonttitle=\bfseries, breakable, #1}
\newtcolorbox{example}[1][]{colback=gray!8, colframe=gray!40, title=Example, fonttitle=\bfseries, breakable, #1}

% Header
\pagestyle{fancy}
\fancyhf{}
\renewcommand{\headrulewidth}{0.4pt}
\lhead{\textbf{Math 104 -- Final Study Guide}}
\rhead{How to Apply the Theorems}
\cfoot{\thepage}

\setlength{\parindent}{0pt}
\setlength{\parskip}{6pt}

\begin{document}

\begin{center}
{\LARGE\bfseries Math 104: Final Study Guide}\\[6pt]
{\large How to Apply the Theorems of Real Analysis}
\end{center}

\tableofcontents
\newpage

%==========================================================================
\section{$\limsup$ and $\liminf$}
%==========================================================================

\begin{intuition}
$\limsup s_n$ is the ``eventual ceiling'' of the sequence --- the largest value that subsequences can cluster around. $\liminf s_n$ is the ``eventual floor.'' If they agree, the sequence converges to that value.

Think of it this way: $\limsup s_n = L$ means that (1) infinitely many terms get close to $L$, and (2) only finitely many terms exceed $L + \eps$ for any $\eps > 0$.
\end{intuition}

\textbf{Two equivalent definitions:}
\begin{enumerate}
\item \textbf{Subsequential limits:} $\limsup s_n = \sup\{$all subsequential limits of $\{s_n\}\}$.
\item \textbf{Tail sups:} $\limsup s_n = \inf_{n} \sup_{k \geq n} s_k = \lim_{n \to \infty} \sup_{k \geq n} s_k$.
\end{enumerate}

The tail-sup formulation is often easier to compute: the sequence $a_n = \sup_{k \geq n} s_k$ is decreasing, so its limit always exists (possibly $\pm\infty$).

\begin{application}
\textbf{When you see $\limsup$ on an exam:}
\begin{itemize}
\item To \textbf{compute} $\limsup s_n$: find all subsequential limits (cluster points), take the largest. Or compute $\sup_{k \geq n} s_k$ and take $n \to \infty$.
\item To \textbf{show convergence}: prove $\limsup s_n = \liminf s_n$.
\item To \textbf{bound a sequence}: $\limsup s_n \leq L$ means $s_n < L + \eps$ eventually.
\item In the \textbf{root/ratio tests}: the test computes $\limsup$ of $|a_n|^{1/n}$ or $|a_{n+1}/a_n|$ and compares it to 1.
\end{itemize}
\end{application}

\begin{example}
$s_n = (-1)^n(1 + 1/n)$. The subsequential limits are $\{-1, 1\}$. So $\limsup s_n = 1$ and $\liminf s_n = -1$.
\end{example}

\begin{example}
$s_n = \sin(n)$. The set of subsequential limits is $[-1, 1]$ (since $\{n \mod 2\pi\}$ is dense in $[0, 2\pi]$). So $\limsup = 1$, $\liminf = -1$.
\end{example}

%==========================================================================
\section{Metric Spaces and Topology}
%==========================================================================

\subsection{Open, Closed, Compact, Connected, Dense}

\begin{intuition}
These five properties describe different ``shapes'' a subset can have:
\begin{itemize}
\item \textbf{Open}: every point has breathing room (an $\eps$-ball fitting inside).
\item \textbf{Closed}: the set ``contains its boundary'' --- limits of sequences in the set stay in the set.
\item \textbf{Compact}: the set is ``finite-like'' --- you can always extract finite subcovers, convergent subsequences, etc. In $\R^n$: closed + bounded.
\item \textbf{Connected}: the set is ``one piece'' --- you can't split it into two separated parts.
\item \textbf{Dense}: the set is ``everywhere'' --- every point of the ambient space is either in the set or arbitrarily close to it.
\end{itemize}
\end{intuition}

\begin{application}
\textbf{To show a set is open:} For each $x \in E$, find $\eps > 0$ with $B(x, \eps) \subseteq E$. Alternatively, show $E$ is a union of open balls, or show $E^c$ is closed.

\textbf{To show a set is closed:} Show $E$ contains all its limit points: if $x_n \in E$ and $x_n \to x$, then $x \in E$. Or show $E^c$ is open.

\textbf{To show a set is compact (in $\R^n$):} Use Heine--Borel: show it's closed AND bounded. For a general metric space, you must work from the definition (every open cover has a finite subcover).

\textbf{To show a set is connected:} In $\R$, show it's an interval (no gaps). In general, assume $E = A \cup B$ with $A, B$ separated and nonempty, derive a contradiction.

\textbf{To show a set is dense:} Show that for every $x \in X$ and every $\eps > 0$, there exists a point of $E$ within distance $\eps$ of $x$. Equivalently, show $\overline{E} = X$.
\end{application}

\begin{pitfall}
\begin{itemize}
\item A set can be \textbf{neither open nor closed} (e.g., $[0, 1)$ in $\R$) or \textbf{both open and closed} ($\emptyset$, $\R$, or any set in the discrete topology).
\item \textbf{Bounded does not imply compact} in general metric spaces. Example: $(0, 1)$ is bounded but not compact. Heine--Borel only works in $\R^n$.
\item \textbf{Closed and bounded} $\neq$ compact in arbitrary metric spaces. You need $\R^n$.
\item An \textbf{infinite intersection of open sets} need not be open: $\bigcap_{n=1}^\infty (-1/n, 1/n) = \{0\}$.
\end{itemize}
\end{pitfall}

\subsection{Closure, Limit Points, and the Closure Formula}

The closure $\overline{E} = E \cup E'$ where $E'$ is the set of limit points. This is the smallest closed set containing $E$.

\begin{application}
\textbf{To find $\overline{E}$:} Ask ``what are the limit points of $E$?'' A point $p$ is a limit point if every neighborhood of $p$ contains a point of $E$ different from $p$.

Example: $E = \{1/n : n \in \N\}$. The only limit point is $0$ (every neighborhood of $0$ contains some $1/n$). So $\overline{E} = E \cup \{0\} = \{0, 1, 1/2, 1/3, \ldots\}$, which is closed.
\end{application}

\subsection{Compactness: Why It Matters}

Compactness is the most powerful tool in analysis. Whenever you have compactness, you get:

\begin{enumerate}
\item Every sequence has a convergent subsequence (in the set)
\item Every continuous function attains its max and min (Extreme Value Theorem)
\item Continuous functions are automatically uniformly continuous
\item Open covers have finite subcovers (the definition!)
\item The finite intersection property holds
\end{enumerate}

\begin{application}
\textbf{The compactness trick:} You need something to hold globally. Get a local version at each point $x$ (using continuity, openness, etc.), which gives an open set $U_x$ around $x$. These $U_x$ cover the compact set $K$. Extract a \textbf{finite} subcover $U_{x_1}, \ldots, U_{x_n}$. Now take $\delta = \min(\delta_{x_1}, \ldots, \delta_{x_n})$, which works globally since the minimum of finitely many positive numbers is positive.

This is exactly how you prove ``continuous on compact $\Rightarrow$ uniformly continuous.''
\end{application}

%==========================================================================
\section{Cantor Set and Its Uses}
%==========================================================================

The Cantor set $C$ is the universal counterexample machine in real analysis.

\textbf{Construction:} Start with $[0, 1]$. Remove the middle third $(1/3, 2/3)$. From each remaining interval, remove its middle third. Repeat. $C = \bigcap_{n=0}^\infty C_n$.

\textbf{Ternary characterization:} $C = \{x \in [0, 1] : x$ has a base-3 expansion using only digits $0$ and $2\}$.

\begin{application}
\textbf{Use the Cantor set when you need:}
\begin{itemize}
\item An \textbf{uncountable set with measure zero}: $C$ is uncountable ($C \sim \{0,1\}^\N$) but has Lebesgue measure $0$ (total length $(2/3)^n \to 0$).
\item A \textbf{perfect set} (closed, no isolated points): $C' = C$.
\item A \textbf{nowhere dense} closed set: $C$ contains no interval, yet is uncountable.
\item A \textbf{totally disconnected} compact set: the only connected subsets of $C$ are singletons.
\item Proof that \textbf{nonempty perfect sets are uncountable}: $C$ is the canonical example.
\end{itemize}
\end{application}

%==========================================================================
\section{Continuity and the Smoothness Hierarchy}
%==========================================================================

\subsection{The Three Faces of Continuity}

Continuity has three equivalent definitions (for metric spaces):
\begin{enumerate}
\item \textbf{$\eps$-$\delta$:} $\forall \eps > 0$, $\exists \delta > 0$: $d(x, y) < \delta \Rightarrow d(f(x), f(y)) < \eps$.
\item \textbf{Sequential:} $x_n \to x \Rightarrow f(x_n) \to f(x)$.
\item \textbf{Topological:} Preimages of open sets are open.
\end{enumerate}

Use $\eps$-$\delta$ for explicit proofs, sequential for contradictions (``if $f$ weren't continuous, there'd be a bad sequence''), and topological for clean abstract arguments.

\subsection{Continuity vs.\ Uniform Continuity}

\begin{intuition}
\textbf{Continuous}: at each point, you can find a $\delta$ that works for that point (different points may need different $\delta$'s).

\textbf{Uniformly continuous}: one $\delta$ works for ALL points simultaneously.
\end{intuition}

\begin{application}
\textbf{To show $f$ is NOT uniformly continuous:} Find two sequences $x_n, y_n$ with $|x_n - y_n| \to 0$ but $|f(x_n) - f(y_n)| \not\to 0$.

Example: $f(x) = x^2$ on $\R$ is not uniformly continuous. Take $x_n = n$, $y_n = n + 1/n$. Then $|x_n - y_n| = 1/n \to 0$ but $|f(x_n) - f(y_n)| = |2 + 1/n^2| \to 2 \neq 0$.

\textbf{Shortcut:} If the domain is compact and $f$ is continuous, you automatically get uniform continuity. So non-uniform-continuity can only happen on non-compact domains.
\end{application}

\subsection{The Smoothness Hierarchy}

$$C^\infty \subsetneq \cdots \subsetneq C^2 \subsetneq C^1 \subsetneq C^0$$

Each inclusion is strict. Know the separating examples:
\begin{itemize}
\item $C^0$ but not $C^1$: $f(x) = |x|$ (continuous, not differentiable at $0$).
\item $C^0$ but not $C^1$ (subtler): $f(x) = x \sin(1/x)$, $f(0) = 0$.
\item Differentiable but derivative not continuous: $f(x) = x^2 \sin(1/x)$, $f(0) = 0$. Here $f'(0) = 0$ exists but $f'$ is not continuous at $0$.
\item $C^\infty$ but not analytic: $f(x) = e^{-1/x^2}$ for $x \neq 0$, $f(0) = 0$.
\end{itemize}

%==========================================================================
\section{Intermediate Value Theorem and Mean Value Theorem}
%==========================================================================

\subsection{IVT: Continuous Functions Hit Everything Between}

\textbf{Statement:} If $f$ is continuous on $[a,b]$ and $c$ is between $f(a)$ and $f(b)$, then $f(p) = c$ for some $p \in (a,b)$.

\begin{application}
\textbf{Use IVT to show equations have solutions:}
\begin{enumerate}
\item Define $f$ so that your equation becomes $f(x) = 0$ (or $f(x) = c$).
\item Find $a, b$ where $f(a) < 0$ and $f(b) > 0$ (or vice versa).
\item Conclude by IVT: $\exists p$ with $f(p) = 0$.
\end{enumerate}

Example: Show $x^5 + x = 1$ has a solution. Let $f(x) = x^5 + x - 1$. $f(0) = -1 < 0$, $f(1) = 1 > 0$. By IVT, $\exists p \in (0,1)$ with $f(p) = 0$.
\end{application}

\begin{pitfall}
IVT requires continuity on a \textbf{closed interval}. It fails for discontinuous functions (Dirichlet function takes values $0$ and $1$ but misses everything in between on any interval).
\end{pitfall}

\subsection{MVT: Derivatives Match Secant Lines}

\textbf{Statement:} $f$ continuous on $[a,b]$, differentiable on $(a,b)$ $\Rightarrow$ $\exists c \in (a,b)$ with $f'(c) = \frac{f(b) - f(a)}{b - a}$.

\begin{application}
\textbf{Use MVT to bound function values:}
$$|f(b) - f(a)| = |f'(c)| \cdot |b - a| \leq \sup_{[a,b]} |f'| \cdot |b - a|.$$

This is the workhorse inequality of analysis. It says: ``if the derivative is small, the function doesn't change much.''

\textbf{Use MVT to prove functions are constant:} If $f' = 0$ on $(a,b)$, then for any $x, y \in (a,b)$: $|f(x) - f(y)| = |f'(c)||x - y| = 0$.

\textbf{Use MVT to prove uniqueness of fixed points:} If $f(x_1) = x_1$ and $f(x_2) = x_2$ with $x_1 \neq x_2$, MVT gives $f'(c) = \frac{f(x_2) - f(x_1)}{x_2 - x_1} = 1$. So if $f' \neq 1$ everywhere, at most one fixed point.
\end{application}

\subsection{Generalized MVT and L'H\^opital's Rule}

The Generalized MVT replaces the denominator $b - a$ with $g(b) - g(a)$:
$$[f(b) - f(a)] g'(c) = [g(b) - g(a)] f'(c).$$

L'H\^opital's rule follows from GMVT: if $f(x), g(x) \to 0$ as $x \to a$, then
$$\frac{f(x)}{g(x)} = \frac{f(x) - f(a)}{g(x) - g(a)} = \frac{f'(c)}{g'(c)} \to \frac{\lim f'}{\lim g'}.$$

\begin{application}
\textbf{L'H\^opital recipe:}
\begin{enumerate}
\item Verify you have $0/0$ or $\infty/\infty$ form.
\item Differentiate numerator and denominator separately.
\item If the new limit exists, it equals the original.
\item If you get $0/0$ again, apply L'H\^opital again.
\end{enumerate}
\end{application}

\subsection{Taylor's Theorem}

$$f(x) = \sum_{k=0}^{n-1} \frac{f^{(k)}(a)}{k!}(x-a)^k + \frac{f^{(n)}(c)}{n!}(x-a)^n$$

The last term is the \textbf{Lagrange remainder} --- it tells you how good the polynomial approximation is.

\begin{application}
\textbf{To bound the error of a Taylor approximation:} The error is $|R_n| = \frac{|f^{(n)}(c)|}{n!}|x-a|^n$ for some unknown $c$. Bound $|f^{(n)}(c)| \leq M$ on the interval, giving $|R_n| \leq \frac{M}{n!}|x-a|^n$.
\end{application}

%==========================================================================
\section{Series}
%==========================================================================

\subsection{Choosing a Convergence Test}

\begin{application}
\textbf{Decision tree for $\sum a_n$:}
\begin{enumerate}
\item Does $a_n \to 0$? If not, the series \textbf{diverges} (necessary condition).
\item Does it look like a known series ($p$-series, geometric)? Use \textbf{comparison}.
\item Are there factorials or exponentials? Try the \textbf{ratio test}.
\item Is $a_n = f(n)$ for some nice $f$? Try the \textbf{root test} (always at least as good as ratio).
\item Is it decreasing and positive? Try \textbf{Cauchy condensation}.
\item Is it alternating? Try the \textbf{alternating series test} (Leibniz).
\end{enumerate}
\end{application}

\subsection{Absolute vs.\ Conditional Convergence}

\textbf{Absolutely convergent:} $\sum |a_n| < \infty$. This is the ``safe'' type --- you can rearrange terms, group them, multiply series, etc.

\textbf{Conditionally convergent:} $\sum a_n$ converges but $\sum |a_n| = \infty$. This is fragile --- the \textbf{Riemann rearrangement theorem} says you can rearrange the terms to make the series converge to \textit{any} value you want (or diverge).

\begin{application}
\textbf{To determine which type:}
\begin{enumerate}
\item First check if $\sum |a_n|$ converges (using comparison, root, ratio tests).
\item If yes: absolutely convergent.
\item If $\sum |a_n|$ diverges but $\sum a_n$ converges (often via alternating series test): conditionally convergent.
\end{enumerate}
\end{application}

\begin{example}
$\sum (-1)^n / n$: converges by alternating series test, but $\sum 1/n$ diverges (harmonic). So it converges \textbf{conditionally}.

$\sum (-1)^n / n^2$: $\sum 1/n^2$ converges ($p$-series, $p = 2 > 1$). So it converges \textbf{absolutely}.
\end{example}

\subsection{Radius of Convergence}

For a power series $\sum c_n (x - a)^n$, the radius of convergence is:
$$R = \frac{1}{\limsup_{n \to \infty} |c_n|^{1/n}}$$

\begin{application}
\textbf{To find $R$:}
\begin{itemize}
\item \textbf{Root test method:} Compute $\limsup |c_n|^{1/n}$, take reciprocal.
\item \textbf{Ratio test shortcut:} If $\lim |c_{n+1}/c_n|$ exists, then $R = \lim |c_n / c_{n+1}|$.
\item \textbf{Always check endpoints} $x = a \pm R$ separately --- convergence there is not determined by $R$.
\end{itemize}
\end{application}

\begin{example}
$\sum x^n / n$: $c_n = 1/n$, so $|c_n|^{1/n} = n^{-1/n} \to 1$, giving $R = 1$. At $x = 1$: $\sum 1/n$ diverges. At $x = -1$: $\sum (-1)^n/n$ converges (alternating). So: converges on $[-1, 1)$.
\end{example}

%==========================================================================
\section{Sequences and Series of Functions}
%==========================================================================

\subsection{Pointwise vs.\ Uniform Convergence}

\begin{intuition}
\textbf{Pointwise:} At each individual point $x$, $f_n(x) \to f(x)$. Different points may converge at very different rates.

\textbf{Uniform:} ALL points converge at the same rate. After some $N$, the entire graph of $f_n$ is within $\eps$ of the graph of $f$.
\end{intuition}

\begin{application}
\textbf{To show uniform convergence:} Compute $\|f_n - f\|_\infty = \sup_{x \in E} |f_n(x) - f(x)|$ and show it $\to 0$.

Common technique: find the maximum of $|f_n(x) - f(x)|$ by taking the derivative and setting it to zero.

\textbf{To show convergence is NOT uniform:}
\begin{enumerate}
\item Find the supremum $\sup_x |f_n(x) - f(x)|$ and show it does NOT go to $0$.
\item Or: the limit function is discontinuous but each $f_n$ is continuous (uniform limit of continuous functions must be continuous).
\end{enumerate}
\end{application}

\begin{example}
$f_n(x) = x^n$ on $[0, 1]$. Pointwise limit: $f(x) = 0$ for $x \in [0,1)$, $f(1) = 1$. The limit is discontinuous, but each $f_n$ is continuous. So the convergence is \textbf{not uniform}.

Alternatively: $\sup_{x \in [0,1]} |x^n - f(x)| \geq |f_n(1-1/n) - 0| = (1 - 1/n)^n \to 1/e \neq 0$.
\end{example}

\subsection{Weierstrass M-Test}

\textbf{Statement:} If $|f_n(x)| \leq M_n$ for all $x$ and $\sum M_n < \infty$, then $\sum f_n$ converges uniformly.

\begin{application}
\textbf{To apply the M-test:}
\begin{enumerate}
\item For each $n$, find $M_n = \sup_{x \in E} |f_n(x)|$.
\item Check if $\sum M_n$ converges.
\item If yes, $\sum f_n$ converges uniformly on $E$.
\end{enumerate}

Example: $\sum_{n=1}^\infty \frac{\sin(nx)}{n^2}$ on $\R$. We have $|\sin(nx)/n^2| \leq 1/n^2 = M_n$. Since $\sum 1/n^2$ converges, the series converges uniformly by M-test.
\end{application}

\begin{pitfall}
The M-test is a \textbf{sufficient} condition, not necessary. The converse is false: a series can converge uniformly even if the M-test fails. But the M-test is the go-to first attempt.
\end{pitfall}

%==========================================================================
\section{Uniform Convergence and Interchange of Limits}
%==========================================================================

The central theme: \textbf{when can we swap the order of two limit operations?} The answer is almost always: when convergence is uniform.

\subsection{Uniform Convergence + Continuity}

\textbf{Theorem:} Uniform limit of continuous functions is continuous.
$$f_n \text{ continuous}, \; f_n \to f \text{ uniformly} \implies f \text{ continuous}$$

\begin{application}
\textbf{Use this to prove a function is continuous:} Write it as the uniform limit of continuous functions (often partial sums of a series). Use the M-test to verify uniform convergence.

\textbf{Use the contrapositive to show convergence is not uniform:} If $f = \lim f_n$ is discontinuous but each $f_n$ is continuous, the convergence cannot be uniform. (This is the $x^n$ argument.)
\end{application}

\subsection{Uniform Convergence + Integration}

\textbf{Theorem:} If $f_n \to f$ uniformly and each $f_n$ is integrable, then $f$ is integrable and
$$\int_a^b f \, dx = \lim_{n \to \infty} \int_a^b f_n \, dx.$$

\begin{intuition}
Uniform convergence lets you ``pass the limit inside the integral.'' This is the easiest interchange theorem --- it just works, no extra hypotheses needed (beyond uniform convergence and integrability of each $f_n$).
\end{intuition}

\subsection{Uniform Convergence + Differentiation}

\textbf{Theorem:} If $f_n$ are differentiable, $f_n'$ converges uniformly to $g$, and $f_n(x_0)$ converges for some $x_0$, then $f_n \to f$ uniformly and $f' = g$.

\begin{pitfall}
This is the trickiest interchange. You need:
\begin{enumerate}
\item Uniform convergence of the \textbf{derivatives} $f_n'$ (not the functions!).
\item Pointwise convergence of $f_n$ at \textbf{at least one point}.
\end{enumerate}

Uniform convergence of $f_n$ alone is NOT enough for $(\lim f_n)' = \lim f_n'$. Example: $f_n(x) = \frac{\sin(nx)}{\sqrt{n}} \to 0$ uniformly, but $f_n'(x) = \sqrt{n} \cos(nx)$ diverges.
\end{pitfall}

\subsection{Summary Table}

\begin{center}
\begin{tabular}{|l|c|c|}
\hline
\textbf{Operation} & \textbf{Hypothesis} & \textbf{Conclusion} \\
\hline
Continuity & $f_n$ cts, $f_n \to f$ unif. & $f$ is cts \\
Integration & $f_n \in \mathscr{R}$, $f_n \to f$ unif. & $\int f = \lim \int f_n$ \\
Differentiation & $f_n' \to g$ unif., $f_n(x_0)$ conv. & $f' = g = \lim f_n'$ \\
\hline
\end{tabular}
\end{center}

%==========================================================================
\section{Equicontinuity and Arzel\`a--Ascoli}
%==========================================================================

\subsection{What is Equicontinuity?}

\begin{intuition}
A family $\{f_n\}$ is equicontinuous if ALL functions in the family are ``uniformly continuous with the same $\delta$.'' In other words, the modulus of continuity is uniform across the entire family, not just across points.

\textbf{Single function, uniform continuity:} One $\delta$ works for all \textbf{points}.\\
\textbf{Family, equicontinuity:} One $\delta$ works for all \textbf{points} AND all \textbf{functions in the family}.
\end{intuition}

\begin{application}
\textbf{To verify equicontinuity:} Find a single $\delta$ (depending only on $\eps$) such that $d(x,y) < \delta \Rightarrow |f_n(x) - f_n(y)| < \eps$ for ALL $n$.

Common sufficient conditions:
\begin{itemize}
\item All $f_n$ are Lipschitz with the \textbf{same} Lipschitz constant $L$: $\delta = \eps/L$.
\item All $f_n$ are differentiable with $|f_n'| \leq M$ (uniformly bounded derivatives): $\delta = \eps/M$ by MVT.
\item $f_n \to f$ uniformly and each $f_n$ is continuous $\Rightarrow$ equicontinuous.
\end{itemize}
\end{application}

\subsection{Arzel\`a--Ascoli: Compactness for Function Spaces}

\textbf{Statement:} If $K$ is compact and $\{f_n\} \subset \mathscr{C}(K)$ is \textbf{pointwise bounded} and \textbf{equicontinuous}, then $\{f_n\}$ has a \textbf{uniformly convergent} subsequence.

\begin{intuition}
Arzel\`a--Ascoli is the Bolzano--Weierstrass theorem for functions. In $\R^n$, bounded sequences have convergent subsequences. In $\mathscr{C}(K)$, pointwise bounded + equicontinuous families have convergent subsequences.

Think of it as: ``bounded + equicontinuous = compact in function space.''
\end{intuition}

\begin{application}
\textbf{To apply Arzel\`a--Ascoli:} Check three things:
\begin{enumerate}
\item \textbf{Domain is compact}: $K$ must be compact.
\item \textbf{Pointwise bounded}: $\forall x \in K$, $\exists M_x$: $|f_n(x)| \leq M_x$ for all $n$.
\item \textbf{Equicontinuous}: One $\delta$ works for all $n$ and all points.
\end{enumerate}
Then conclude: there exists a uniformly convergent subsequence.
\end{application}

\begin{example}
Show that $\{f_n\}$ where $f_n: [0,1] \to \R$, $|f_n| \leq 10$, and $|f_n'| \leq 5$ for all $n$ has a uniformly convergent subsequence.

\textbf{Solution:} Domain $[0,1]$ is compact. Pointwise bounded: $|f_n(x)| \leq 10$. Equicontinuous: $|f_n(x) - f_n(y)| \leq 5|x - y|$ by MVT, so $\delta = \eps/5$ works for all $n$. By Arzel\`a--Ascoli, done.
\end{example}

%==========================================================================
\section{Riemann Integrability}
%==========================================================================

\subsection{The Big Picture}

A bounded function $f$ on $[a,b]$ is Riemann integrable if and only if it can be ``squeezed'' between step functions that are arbitrarily close together.

\subsection{The Criterion}

$f \in \mathscr{R}[a,b]$ $\Leftrightarrow$ $\forall \eps > 0$, $\exists$ partition $P$ with $U(P,f) - L(P,f) < \eps$.

\begin{application}
\textbf{To show $f \in \mathscr{R}$:} Construct a partition that makes $U - L$ small. Strategy:
\begin{enumerate}
\item Identify where $f$ has big oscillation (``bad'' intervals).
\item Make the total length of bad intervals very small.
\item On the ``good'' intervals, $f$ doesn't vary much, so $M_i - m_i$ is small.
\item Total: $U - L \leq (\text{bad contribution}) + (\text{good contribution})$.
\item Bad: $\leq 2\|f\|_\infty \cdot (\text{total bad length})$. Good: $\leq \eps \cdot (b-a)$.
\end{enumerate}

\textbf{To show $f \notin \mathscr{R}$:} Show $U(P,f) - L(P,f) \geq c > 0$ for ALL partitions $P$.
\end{application}

\subsection{Sufficient Conditions (know these!)}

\begin{enumerate}
\item $f$ continuous on $[a,b]$ $\Rightarrow$ $f \in \mathscr{R}$. (Uniform continuity controls oscillation.)
\item $f$ monotone on $[a,b]$ $\Rightarrow$ $f \in \mathscr{R}$. (Oscillation on $[x_{i-1}, x_i]$ is $f(x_i) - f(x_{i-1})$, and these telescope.)
\item $f$ bounded with discontinuities forming a \textbf{zero set} $\Rightarrow$ $f \in \mathscr{R}$.
\end{enumerate}

A set $E$ is a \textbf{zero set} (measure zero) if for every $\eps > 0$, $E$ can be covered by countably many intervals with total length $< \eps$. Examples: finite sets, countable sets (like $\Q$), the Cantor set.

\subsection{The Three Key Counterexamples}

\begin{enumerate}
\item \textbf{Thomae's function} $T(x)$: $T(p/q) = 1/q$ (in lowest terms), $T(\text{irrational}) = 0$.
\begin{itemize}
\item Continuous at every irrational (given $\eps$, only finitely many rationals $p/q$ with $1/q \geq \eps$ in any bounded interval).
\item Discontinuous at every rational.
\item $\in \mathscr{R}$ because the discontinuity set $\Q$ is a zero set.
\item Shows: you can be discontinuous on a dense set and still be integrable.
\end{itemize}

\item \textbf{Dirichlet function} $D(x) = \mathbf{1}_\Q(x)$:
\begin{itemize}
\item Discontinuous everywhere.
\item $U(P, D) = 1$ and $L(P, D) = 0$ for every partition (every interval contains both rationals and irrationals).
\item $\notin \mathscr{R}$.
\item Shows: the ``everywhere discontinuous'' case fails.
\end{itemize}

\item $\mathbf{1/\sqrt{x}}$ \textbf{on} $(0, 1]$:
\begin{itemize}
\item Unbounded near $0$.
\item Not Riemann integrable (Riemann integration requires bounded functions on closed intervals).
\item But the improper integral $\int_0^1 x^{-1/2} dx = 2$ exists.
\item Shows: unboundedness breaks Riemann integrability, but improper integrals can rescue the situation.
\end{itemize}
\end{enumerate}

\begin{application}
\textbf{Geometric interpretation of $U(P,f) - L(P,f) < \eps$:} Draw the function. The upper sum is the area of rectangles sitting on top of the curve; the lower sum sits below. Their difference is the total area of the ``cracks'' between upper and lower rectangles. Integrability means these cracks can be made arbitrarily thin.
\end{application}

%==========================================================================
\section{Putting It All Together: Common Exam Problem Types}
%==========================================================================

\begin{enumerate}
\item \textbf{``Show $f$ is uniformly continuous''}: If the domain is compact, just show continuity. Otherwise, find $\delta$ independent of the point.

\item \textbf{``Show $\sum f_n$ converges uniformly''}: Try the Weierstrass M-test first.

\item \textbf{``Show the limit is continuous/integrable/differentiable''}: Show uniform convergence (or for differentiation: uniform convergence of derivatives), then apply the interchange theorem.

\item \textbf{``Does the family have a convergent subsequence?''}: Check Arzel\`a--Ascoli conditions (compact domain, pointwise bounded, equicontinuous).

\item \textbf{``Show $f$ is Riemann integrable''}: Check if $f$ is continuous, monotone, or has a zero set of discontinuities.

\item \textbf{``Find the radius of convergence''}: Use $R = 1/\limsup|c_n|^{1/n}$ or the ratio test. Then check endpoints.

\item \textbf{``Show a sequence converges''}: Show $\limsup = \liminf$, or show it's Cauchy, or show it's monotone and bounded.

\item \textbf{``Show a set is compact/connected/dense''}: For compact in $\R^n$: closed and bounded. For connected in $\R$: show it's an interval. For dense: show $\overline{E} = X$.
\end{enumerate}

\end{document}
